\documentclass[12pt]{article}
\usepackage[a4paper,margin=1in]{geometry}
\usepackage{titlesec}
\usepackage{hyperref}
\titleformat{\section}{\large\bfseries}{\thesection}{1em}{}

\title{Book Proposal: Applying TRANSFORMATICS in GENETICS}
\author{Joseph Willrich Lutalo}
\date{}

\begin{document}

\maketitle

\section{Author Information}
\textbf{Joseph Willrich Lutalo} \\
Independent researcher and theorist in software language engineering, stochastic processes and mathematical statistics. As of this moment, he is based at a private, personal laboratory known as ``Nuchwezi Research" in Garuga, UGANDA, but otherwise was recently (2024) a finalist graduate researcher at Makerere University's College of Computing and Information Sciences (CoCIS), from where he completed his MSc. Software Engineering and Data Communications. His Master's thesis, ``VOSA: A Reusable and Reconfigurable Voice Operated Support Assistant Chatbot Platform", introduced a hands-free human support tool driven by user-configurable knowledge models based on QAKBs. Since founding Nuchwezi (\url{nuchwezi.com}), a research-driven technology service provider based in UGANDA, circa 2014, he has made extensive contributions in the form of mobile and web applications, software libraries, a new general-purpose programming language and several other contributions in both the arts and science. He has a total of 50+ academic works (ORCID), 18 scholarly articles (Google Scholar) and has contributed a total of 100+ mini-lectures on several topics in computer science, maths and software engineering (YouTube). Joseph has written 2 full-length novels as of this moment, and 3 novellas. This particular book on Transformatics in Genetics is his first full-length book focused on science and mathematics.\\

Refer to ORCID: \url{https://orcid.org/0000-0002-0002-4657}


\section{Synopsis}
This mini-treatise introduces and applies the emerging mathematical theory of \textbf{Transformatics}---a discipline focused on the study, analysis and application of symbolic and quantitative transformations of [especially] ordered sequences via abstract machines known as sequence transformers, with focus on its application in the domain of \textbf{genetics}. By bridging symbolic computation with biological sequence analysis, the work opens doors to new ways to approach the understanding, modeling, and manipulation of genetic code (as well as any arbitrary information encoded using ordered sequences in general).\\

The treatise presents:
\begin{itemize}
    \item A rigorous formalism for understanding and applying sequence transformation using transformers (\textbf{16 transformers presented in total}), expressed and specified in mathematics and not just software languages.
    \item \textbf{16 Working Definitions} in relation to various concepts in analytical genetics, sequence analysis and sequence processing in general (including 3 extra definitions not directly under Genetics, but also meant to demonstrate how Transformatics could be leveraged/applied in Physics).
    \item \textbf{4 Theorems}  developed and presented for the first time, spanning properties of genetic sequences and then sequence complements in general.
    \item \textbf{3 Foundational laws} governing na-Sequences and genome sequencing.
    \item \textbf{30 Transformations} presented to help illustrate, but also demonstrate particular computational solutions to many kinds of basic problems in sequence analysis, processing and several concerning gene/genome expression problems.
    \item \textbf{8 Algorithms} developed and presented for the first time including solutions in sequence alignment, pseudo-random sequence generation, encoding and simulation of gene expression, and more.
    \item \textbf{3 Tabular Analyses} demonstrating how to leverage statistical analysis and the concepts of transformatics in quantitative analysis of genetic sequences among other problems.
    \item \textbf{3 Core Problems} in the domain of \textbf{Genome Sequencing} are formally defined and their theoretical solutions also presented for the first time.
    \item \textbf{More than 10 Source Code Examples}, some in the TEA language (by the author), and relating to conducting quick analyses of arbitrary na-Sequences via the commandline, and 8 related to how to generate arbitrary random genetic-code sequences using Python.
    \item As well as \textbf{9 important future problems} that have been identified and highlighted as future work, spanning pathological gene programs, automated genetic sequence analysis, use of Godel numbers in optimal genetics databases, a computational model of the ribosome and more.
\end{itemize}

This work is both theoretical and practical, designed to inspire researchers in genetics, bioinformatics, and computational biology to adopt transformatics as a tool for discovery and innovation.

\section{Target Audience}
\begin{itemize}
    \item Researchers in \textbf{genetics}, \textbf{bioinformatics}, and \textbf{computational biology}
    \item Scholars in \textbf{automata theory}, \textbf{symbolic computation}, and \textbf{mathematical modeling}
    \item Graduate students and professionals seeking interdisciplinary approaches to genome analysis
    \item \textbf{Software Engineers} and \textbf{Computer Scientists} looking for problems and opportunities in genetics, bioinformatics and generally in sequence analysis and processing.
    \item \textbf{Mathematicians} interested in general problems around sequence processing and analysis.
\end{itemize}

\section{Market Potential}
The book fills a gap between symbolic computation and biological sequence analysis. It offers:
\begin{itemize}
    \item A new lens for genome sequencing and gene expression modeling
    \item Tools for algorithmic exploration of DNA and genetic structures
    \item A theoretical foundation for future applications in genetic engineering and synthetic biology
    \item A new mathematics --- not just notations, but also a calculus/mechanics, with which to tackle, think about or model new and old problems, as well as their solutions, in genetics, computing and the general domain of sequence processing and analysis.
\end{itemize}

\section{Table of Contents (Tentative)}
\begin{enumerate}
    \item Introduction (to Genetics and Transfomatics for Non-Experts)
    \item Sequence Symbol Sets and Sequence Transforms Applied to Genetic Code
    \item Sequence Abstraction using Symbol Sets: From Flat-Structure Sequences to Higher-Order
Sequences
    \item Sequence Analysis Using Symbol Sets
    \item Sequence Analysis Using Complement Sets
    \item Sequence Analysis Using The Anagram Distance and the Modal Sequence Statistic
    \item The Mathematics of Genome Sequencing: Sequence Alignment, The Modal Sequence, The
Sequencing Machine and the Identifier Genome Sequence Law
    \item Random Sequence Generators (RSGs)
    \item Gene Expression in Living Organisms Leveraging Genetic Code (DNA → mRNA → Protein →
Organism)
    \item Transformatics and Lu-Genome Expression System in Bio-Automata
    \item Conclusion, Future Directions and Open Problems
    \item Appendices: Genetics, Physics and Transformatics
\end{enumerate}

\section{Sample Chapters}
Available upon request, however, a complete draft of the proposed book is also available in electronic book format via I*POW (author's own staging platform for traditionally unpublished work --- see \url{https://doi.org/10.6084/m9.figshare.29941127}). 

\section{Length and Timeline}
Estimated manuscript length: \textbf{200–300 pages}, but currently: 206 pages, and aprox. 56000 words (first edtn draft)\\
Manuscript status: \textbf{Completed and ready for review}

\section{Why This Book Matters}
This treatise proposes a paradigm shift in how genetic sequences are explored theoretically and especially mathematically. It offers a rich theoretical foundation and practical tools that could redefine computational genetics and inspire new research directions.

\end{document}